\section{总体分析}

本题的核心是在多约束条件下求多年种植结构的最优决策，三个问题层层递进，不确定性逐步增强，本文按以下思路展开：

\begin{itemize}
	\item \textbf{数据视角}：附件数据存在缺失与合并单元格，需先清洗。以2023年实际种植量为各作物预期销售量的基准；以2023年产量、成本、价格作为确定性口径的基准值，并据此推算各作物收益结构，为建模提供参数。
	\item \textbf{问题1（确定性规划）}：把54块地、41种作物、7年、两季次的安排抽象为0-1整数规划（MILP），逐条落实土地容量、季节适应性、水浇地二季复种、重茬回避、豆类三年轮作与销售上限等约束，目标为7年总净收益最大，分别求解滞销浪费与折价出售两种情形。
	\item \textbf{问题2（随机规划）}：销售量、亩产量、价格、成本均存在不确定性。采用情景模拟把不确定性离散为若干随机情景，建立两阶段随机规划模型——第一阶段确定种植方案，第二阶段依据情景实现确定销售与收益；以"期望收益$-\lambda\cdot$CVaR$_\alpha$(损失)"为目标，兼顾期望与尾部风险，通过调节风险系数$\lambda$得到不同风险偏好下的最优方案。
	\item \textbf{问题3（相关结构）}：作物之间并非相互独立。从食用消费习惯（替代、互补）、价格联动、产销量联动等角度构造相关性矩阵，用Cholesky分解生成相关情景，比较相关与独立两种假设下最优方案与风险指标的差异，论证考虑相关性的合理性。
	\item \textbf{模型评价}：对各问方案做约束校验、灵敏度与稳健性分析，并讨论模型适用性与推广方向。
\end{itemize}
